% Source: https://artofproblemsolving.com/community/c3h1632213_1231000
% Stable identity: aops-topic-1632213-q01
Going $12$ places at a time, the $k$-th mark ($k=0,1,2,\dots$) falls on the number congruent to $1+12k$ modulo $1000$.

Two marks $j<k$ land on the same number exactly when $12k\equiv12j\pmod{1000}$, i.e. $1000\mid12(k-j)$, i.e. $250\mid 3(k-j)$, i.e. $250\mid k-j$. So the marks for $k=0,1,\dots,249$ are all different, and the first repeat is at $k=250$, when the process returns to $1$. Hence exactly $250$ numbers are marked.

(These are exactly the numbers congruent to $1$ modulo $4$: since $1000$ and $12$ are multiples of $4$, every mark is $\equiv1\pmod4$, and there are exactly $250$ such numbers from $1$ to $1000$.)

So $1000-250=\boxed{750}$ numbers are never marked.
